% Generated draft: requires mathematical and source review.
\documentclass[11pt]{article}
\usepackage[margin=25mm]{geometry}
\usepackage{fontspec}
\setmainfont{Latin Modern Roman}
\usepackage{amsmath,amssymb,mathtools,unicode-math}
\setmathfont{Latin Modern Math}
\usepackage{enumitem,needspace,xurl,hyperref}
\setlength{\parindent}{0pt}
\setlength{\parskip}{5pt}
\setlength{\emergencystretch}{3em}
\newcommand{\N}{\mathbb{N}}
\newcommand{\Z}{\mathbb{Z}}
\newcommand{\Q}{\mathbb{Q}}
\newcommand{\R}{\mathbb{R}}
\newcommand{\C}{\mathbb{C}}
\newcommand{\F}{\mathbb{F}}
\newcommand{\A}{\mathbb{A}}
\newcommand{\PP}{\mathbb{P}}
\newcommand{\E}{\mathbb{E}}
\newcommand{\eps}{\varepsilon}
\newcommand{\dd}{\,\mathrm{d}}
\begin{document}
\section*{UnsolvedMath ID: 800114}
\textbf{Source problem code:} catalogue entry
\subsection{Definitions and mathematical statement}
Write \(\mathbb T^2=\R^2/\Z^2\). For \(\alpha\in\mathbb T^2\), the rotation \(R_\alpha\) is the translation \(R_\alpha(x)=x+\alpha\). For a compact metrizable surface \(S\), let \(\operatorname{Homeo}(S)\) denote its homeomorphism group. Given a compatible metric \(d_S\), define the uniform distance by
\[
d_{C^0}(u,v)=\sup_{x\in S}d_S\bigl(u(x),v(x)\bigr).
\]
For a specified subgroup \(H(S)\leq\operatorname{Homeo}(S)\), define the conjugacy orbit
\[
\mathcal C_{H(S)}(f)=\{h\circ f\circ h^{-1}:h\in H(S)\}.
\]
Its closure below is taken in the uniform topology on \(\operatorname{Homeo}(S)\).

The supplied record does not define ``non-spreading.'' Write \(\mathcal N(S)\) for the class intended by that term, without assigning it an unsupported definition. The record also does not specify the conjugating subgroup or the topology of approximation. Conditional on interpreting approximation in the uniform topology and taking \(H(S)\) to be the intended conjugating subgroup, the two questions are:
\begin{enumerate}
\item Is it true that
\[
\forall f\in\mathcal N(\mathbb T^2)\quad
\exists\alpha\in\mathbb T^2\quad
\forall\varepsilon>0\quad
\exists h\in H(\mathbb T^2):\quad
 d_{C^0}(h\circ f\circ h^{-1},R_\alpha)<\varepsilon?
\]
Equivalently, must the closure of the conjugacy orbit contain a rotation?
\item For every higher genus surface \(S\) intended by the source, is it true that
\[
\forall f\in\mathcal N(S)\quad
\forall\varepsilon>0\quad
\exists h\in H(S):\quad
 d_{C^0}(h\circ f\circ h^{-1},\operatorname{id}_S)<\varepsilon?
\]
Equivalently, must \(\operatorname{id}_S\) belong to the closure of the conjugacy orbit?
\end{enumerate}
Taking \(H(S)=\operatorname{Homeo}(S)\) gives the interpretation using all homeomorphisms as conjugators. If ``higher genus surface'' means a closed connected orientable surface, the second question ranges over every genus \(g\geq2\). These choices are not explicit in the supplied record. The undefined non-spreading condition and the unspecified surface and conjugacy conventions prevent a fully determined formulation from this record alone.
\subsection{Short English statement}
The questions ask whether every non-spreading torus homeomorphism can be conjugated arbitrarily close to a translation, and whether every non-spreading homeomorphism of a higher genus surface can be conjugated arbitrarily close to the identity. The record leaves the definition of non-spreading and several ambient conventions unspecified. Its \emph{open} label is supplied registry metadata, not an independently verified status assessment.
\subsection{Sources}
\begin{itemize}
\item Local UnsolvedMath registry, permanent ID \(800114\), exact source problem code \texttt{catalogue entry}, title ``Foliations and Diffeomorphism Groups --- Question 55.'' The supplied record attributes the problem to Nariman, Yazdi, with proposed year 2025, and labels it open.
\item Nariman, Yazdi, \emph{Problem List on Foliations and Diffeomorphism Groups} (2025), Question 55. Supplied source URL: \url{https://arxiv.org/abs/2503.18274}. The supplied provenance records access on 2026-10-16 and extraction by \texttt{arxiv\_latex}; these are supplied provenance statements, not a claim of independent access or verification.
\end{itemize}
\end{document}
